% figures/runtime/schwerpunkt_concept-DE.tex

\begin{figure}[ht]
\centering
\begin{tikzpicture}[
  node distance=0.9cm,
  box/.style={
    draw,
    rounded corners,
    align=center,
    minimum width=5.4cm,
    minimum height=0.75cm
  },
  arrow/.style={->, thick}
]
\node[box] (curv) {Krümmung + Überhöhung\\$\kappa(s), u(s)$};
\node[box, below=of curv] (pose3) {Realisierter Laufzeitzustand\\$\mathrm{pose3}_{\mathrm{real}}$};
\node[box, below=of pose3] (cg) {Schwerpunkttrajektorie\\$x_{\mathrm{CG}}(s)$};
\node[box, below=of cg] (balance) {Gleichgewichtstrajektorie\\$a_{\mathrm{eff}}(s)$};
\node[box, below=of balance] (comfort) {Komfort / Zulässigkeit};

\draw[arrow] (curv) -- (pose3);
\draw[arrow] (pose3) -- (cg);
\draw[arrow] (cg) -- (balance);
\draw[arrow] (balance) -- (comfort);
\end{tikzpicture}

\caption{Schwerpunktinterpretation: Trassierungsentwurf wird zur Formung operativer Zustandstrajektorien.}
\label{fig:schwerpunkt-concept}
\end{figure}
